\documentclass[11pt]{article}

\usepackage[utf8]{inputenc}
\usepackage[T1]{fontenc}
\usepackage[margin=1in]{geometry}
\usepackage{amsmath}
\usepackage{amssymb}
\usepackage{booktabs}
\usepackage{enumitem}
\usepackage{float}
\usepackage{xcolor}
\usepackage{hyperref}

\hypersetup{
  colorlinks=true,
  linkcolor=blue!50!black,
  citecolor=blue!50!black,
  urlcolor=blue!60!black,
  pdftitle={PhiGraph Core 4.0: A Shadow-First Evidence Ledger and Policy-Gated Runtime for Software-Agent Operations},
  pdfauthor={Walter Calmels von Dem Knesebeck}
}

\title{\textbf{PhiGraph Core 4.0: A Shadow-First Evidence Ledger and Policy-Gated Runtime for Software-Agent Operations}}
\author{
  Walter Calmels von Dem Knesebeck\\
  TUCH Systems, Santiago, Chile\\
  \href{mailto:wcalmels@phi47.cl}{wcalmels@phi47.cl}
}
\date{July 2026}

\begin{document}
\maketitle

\begin{abstract}
Software agents and artificial-intelligence systems can produce claims,
recommendations, and action proposals faster than organizations can verify and
govern them. This paper presents PhiGraph Core 4.0.0, an MIT-licensed,
model-agnostic system that treats model output as a proposal rather than as
verified truth. PhiGraph represents claims, evidence, verifications, action
proposals, policy decisions, and outcomes as linked records; stores them in a
scope-tagged, tamper-evident ledger; and places policy evaluation between agent
proposals and operational execution. Replay and shadow modes cannot execute
external actions, while higher-authority modes require both an explicit
executor and an allowing policy decision. We describe the protocol, runtime,
ledger, deployment boundaries, and Python implementation. Engineering
verification comprises 117 automated tests; the continuous-integration
configuration targets Python 3.10--3.13. An initial external experiment on the CIC-IDS2017
Friday DDoS flow file compares a PhiGraph relational score with three
unsupervised baselines. Local Outlier Factor is strongest on this tabular
benchmark; PhiGraph ranks second by PR-AUC and attains 0.840 precision at a
10\% alert budget. The experiment demonstrates functional integration, not
state-of-the-art superiority: it uses a single imperfect dataset, a
feature-similarity graph, and one fixed split. PhiGraph is therefore presented
as an auditable governance substrate and research prototype, not as a proof of
correctness, causal identification, or safe autonomous operation.
\end{abstract}

\noindent\textbf{Keywords:} AI governance; software agents; provenance;
evidence ledger; graph intelligence; policy enforcement; shadow deployment;
reproducible software

\section{Introduction}

Modern software agents can synthesize text, modify code, recommend operational
changes, invoke tools, and coordinate multi-step workflows. Their outputs may
be useful without being reliable, authorized, or supported by reproducible
evidence. This creates a control problem: an organization must determine what
was claimed, which evidence supported or refuted it, which policy governed a
proposed action, who approved it, and what happened after evaluation or
execution.

Existing guidance emphasizes lifecycle risk management and traceability.
The NIST AI Risk Management Framework organizes activities around governing,
mapping, measuring, and managing AI risk \cite{nist2023airmf}. W3C PROV provides
a general vocabulary for provenance \cite{lebo2013provo}, while model cards and
dataset datasheets improve documentation of released artifacts
\cite{mitchell2019modelcards,gebru2021datasheets}. These approaches are
important, but they do not by themselves define an operational boundary
between a model-generated proposal and an authorized action.

PhiGraph addresses this boundary through an evidence-oriented protocol and a
policy-gated runtime. Its central invariant is that model output is a candidate
claim or action proposal, not verified truth or implicit authority. The system
records relationships among claims, evidence, verification results, policies,
approvals, and outcomes so that a decision can be reconstructed after the
fact.

This paper makes four contributions:

\begin{enumerate}[leftmargin=*]
  \item a typed protocol for claims, evidence, verifications, action proposals,
        policy decisions, and outcomes;
  \item a runtime in which policy evaluation and an explicit executor are
        necessary conditions for external action;
  \item a backend-neutral, scope-tagged ledger with hash-chain integrity
        metadata and optional HMAC evidence signatures; and
  \item a bounded empirical evaluation comprising software verification and an
        external anomaly-ranking experiment, together with explicit threats to
        validity.
\end{enumerate}

\section{Related Work}

\subsection{Governance, documentation, and provenance}

The NIST AI RMF is a voluntary, use-case-agnostic framework for managing AI
risk \cite{nist2023airmf}. PhiGraph is not an implementation or certification
of that framework; it provides operational records that may support governance
activities. PROV-O formalizes entities, activities, agents, and provenance
relations \cite{lebo2013provo}. PhiGraph uses a narrower application protocol
centered on verification and policy-gated action. A future interoperability
layer could map PhiGraph records to PROV-O.

Model cards \cite{mitchell2019modelcards} and datasheets for datasets
\cite{gebru2021datasheets} make intended uses, evaluation conditions, and
limitations explicit. PhiGraph applies a related transparency principle at
runtime: each claim and action remains linked to its evidence, decision, and
outcome rather than relying only on static release documentation.

\subsection{Relational representations and software assurance}

Graph-based representations make entities and relationships first-class and
support structured reasoning \cite{battaglia2018relational}. PhiGraph uses
heterogeneous graphs and classical graph analytics; it does not require a graph
neural network and does not claim that graph structure alone establishes
causality. In software assurance, in-toto protects software-supply-chain
integrity through verifiable step metadata \cite{torresarias2019intoto}.
PhiGraph shares an emphasis on attributable records but targets a broader
claim--evidence--policy--action lifecycle.

Agentic systems introduce risks including tool misuse, privilege compromise,
untraceable actions, and cascading errors. The OWASP Agentic Security
Initiative provides an industry threat taxonomy \cite{owasp2025agentic}.
PhiGraph's deny-by-default policy, approval requirements, idempotency controls,
and shadow-first deployment are engineering controls relevant to these risks;
they are not a complete security solution.

\section{System Model and Design Goals}

Let an agent produce a proposal
\[
P = (C, A),
\]
where $C$ is a set of claims and $A$ is a set of proposed actions. A claim does
not become verified merely because an agent assigns it high confidence.
Instead, a verifier records a result over a claim and an explicit evidence set.
Similarly, an action is not authorized by its presence in $P$: the policy
engine evaluates the action under a runtime mode and an approval set.

PhiGraph is designed around six goals:

\begin{enumerate}[leftmargin=*]
  \item \textbf{Evidence boundedness:} when evidence identifiers are supplied,
        verification results reference evidence already registered in the
        ledger.
  \item \textbf{Least authority:} absence of a matching policy blocks an
        action.
  \item \textbf{Traceability:} claims, decisions, and outcomes remain
        reconstructable.
  \item \textbf{Safe evaluation:} replay and shadow runs never execute an
        external action.
  \item \textbf{Scope tagging:} tenant and project identifiers propagate
        through stored records and can be used to filter queries.
  \item \textbf{Model neutrality:} deterministic programs, statistical models,
        and language-model agents can use the same protocol.
\end{enumerate}

\section{PhiGraph Architecture}

\begin{figure}[t]
\centering
\fbox{\begin{minipage}{0.91\linewidth}
\centering
\textbf{Agent or deterministic adapter}\\
$\downarrow$ proposal: claims and actions\\[2pt]
\textbf{Protocol layer}\\
Claim $\leftrightarrow$ Evidence $\leftrightarrow$ Verification\\
ActionProposal $\leftrightarrow$ PolicyDecision $\leftrightarrow$ Outcome\\[2pt]
$\downarrow$\\
\textbf{Core service}\\
Evidence ledger \quad Policy engine \quad Runtime \quad Telemetry\\[2pt]
$\downarrow$\\
\textbf{Deployment boundary}\\
Replay / Shadow (non-executing) \quad Copilot / Guarded-auto (explicit executor)
\end{minipage}}
\caption{PhiGraph separates agent proposals from verification, policy
evaluation, and execution. Arrows indicate recorded processing relationships,
not automatic trust propagation.}
\label{fig:architecture}
\end{figure}

\subsection{Canonical protocol}

Protocol version 2.0.0 exposes six principal frozen dataclass record schemas:

\begin{itemize}[leftmargin=*]
  \item \texttt{Claim}: statement, type, subject, issuer, status, confidence,
        evidence references, and optional supersession;
  \item \texttt{Evidence}: kind, source, payload, status, content hash,
        observation time, and metadata;
  \item \texttt{Verification}: claim, verifier, method, result, evidence
        references, rationale, and time;
  \item \texttt{ActionProposal}: action type, target, proposer, parameters,
        rationale claims, reversibility, and risk level;
  \item \texttt{PolicyDecision}: effect, applicable policies, reasons,
        required approvals, and time; and
  \item \texttt{Outcome}: action, status, whether execution occurred, observed
        effects, evidence references, and time.
\end{itemize}

Claim states distinguish proposed, unverified, partially verified, verified,
refuted, and superseded records. Policy effects distinguish allow, warn,
require approval, and block. These states make uncertainty and governance
decisions explicit instead of collapsing them into a single confidence score.
The current implementation does not require a non-empty evidence set for every
verified claim and does not assess evidentiary sufficiency. This is an
enforcement gap, not a guarantee supplied by the protocol.

\subsection{Evidence ledger}

The ledger supports JSON, SQLite, and PostgreSQL backends through a common
interface. Collections are maintained for claims, evidence, verifications,
actions, policy decisions, and outcomes. Most records are appended. Claim
verification performs a controlled in-place state transition and then rebuilds
the affected collection chains. Each ordinary append computes
\[
h_i =
\operatorname{SHA256}\!\left(h_{i-1}\,\Vert\,k\,\Vert\,
\operatorname{canonical}(r_i)\right),
\]
where $r_i$ is a record and $k$ identifies its collection. Verification checks
link continuity and recomputes each hash. Evidence payloads receive a SHA-256
content hash and may be signed with HMAC-SHA256 when a signing key is supplied.

Chain verification detects changes relative to the currently stored chain.
This mechanism is \emph{tamper evident}, not tamper proof: without an external
anchor or protected signing key, it cannot distinguish malicious tampering
from an authorized rebuild. An attacker with write access to both the records
and all integrity keys can rebuild a consistent chain. Stronger deployment
requires external key management, append-only storage, access control, backups,
and independent audit anchors.

\subsection{Policy-gated runtime}

The policy engine matches action type, permitted runtime modes, maximum risk,
reversibility, and required approvals. No matching rule produces an implicit
deny decision. For an action $a$, decision $d$, mode $m$, and executor $x$, the
core execution condition is
\[
\operatorname{execute}(a) \iff
\bigl(d.\mathrm{effect}=\mathrm{allow}\bigr)
\land \bigl(m\in\{\mathrm{copilot},\mathrm{guarded\_auto}\}\bigr)
\land (x\neq\varnothing).
\]
Replay and shadow modes therefore generate simulated outcomes with
\texttt{executed=false}. The packaged deployment configuration is stricter:
it defaults to shadow-only operation, rejects attempts to disable that
boundary, and disables real connectors. The core abstractions expose
higher-authority modes for controlled future integration, but the supplied
deployment does not grant arbitrary external execution.

\subsection{API, SDK, and operations}

The Python SDK can send tenant and project headers, API-key or bearer
credentials, and idempotency keys. Authentication, PostgreSQL persistence, and
signing are supported by lower-level Core components when explicitly
configured. The packaged application does not currently wire every one of
these options into the Core v3 router. Its Docker configuration must therefore
be treated as a local evaluation setup unless protected by an external
authentication and network boundary. Scope headers are application metadata
and query filters; the present implementation does not provide end-to-end
tenant isolation or same-scope referential-integrity guarantees.

\section{Evaluation}

\subsection{Research questions}

The evaluation addresses two bounded questions:

\begin{description}[leftmargin=2.3cm,style=nextline]
  \item[RQ1] Is the software version mechanically testable across its supported Python
  versions and packaging paths?
  \item[RQ2] Does the current relational anomaly score produce useful rankings
  on one external network-flow benchmark relative to simple unsupervised
  baselines?
\end{description}

It does not test whether PhiGraph prevents all unsafe agent actions, proves
causality, or outperforms specialized state-of-the-art intrusion-detection
systems.

\subsection{Software verification}

The software tree contains 117 automated tests covering protocol records, ledger
integrity, policy evaluation, runtime modes, persistence backends, API and SDK
boundaries, agent workflows, graph operations, benchmarks, and deployment
configuration. Continuous integration is configured to install the declared
extras and run the suite on Python 3.10, 3.11, 3.12, and 3.13; separate jobs
build the Python distribution, install its wheel in an isolated environment,
and build the Docker image. A local run on July 27, 2026 executed all 117 tests
successfully. A separate synthetic benchmark attempt under local Python 3.14,
which is outside the CI matrix, stopped with an ARPACK eigensolver error and is
excluded from the reported empirical results. Passing tests establish
conformance to encoded cases, not correctness for all deployments.

\subsection{External anomaly-ranking experiment}

The committed external-validation artifact uses the CIC-IDS2017 Friday
afternoon DDoS flow file \cite{sharafaldin2018cicids}. It contains 225,745
rows: 97,718 benign and 128,027 DDoS. With random seed 47, the protocol draws a
12,000-row benign-only reference set and evaluates a balanced test set of
6,000 benign and 6,000 DDoS flows. Preprocessing and variable-feature selection
use only the benign reference. Methods are ranked continuously and evaluated
at a top-10\% alert budget.

The compared methods are robust multivariate deviation, Isolation Forest,
Local Outlier Factor (LOF), and a PhiGraph relational score combining robust
deviation, mean nearest-neighbor distance, and disagreement between their
percentile ranks. Table
\ref{tab:cicids} reproduces the committed result.

\begin{table}[H]
\centering
\small
\caption{CIC-IDS2017 Friday DDoS anomaly-ranking results. Runtime is reported
by the committed validation artifact; its execution environment was not
recorded.}
\label{tab:cicids}
\begin{tabular}{lrrrrrr}
\toprule
Method & ROC-AUC & PR-AUC & P@10\% & R@10\% & F1@10\% & Time (s)\\
\midrule
Local Outlier Factor & 0.9757 & 0.9506 & 0.9617 & 0.1923 & 0.3206 & 0.322\\
PhiGraph relational  & 0.7295 & 0.7316 & 0.8400 & 0.1680 & 0.2800 & 0.165\\
Isolation Forest     & 0.7403 & 0.6840 & 0.7817 & 0.1563 & 0.2606 & 0.321\\
Robust deviation     & 0.5771 & 0.5540 & 0.4117 & 0.0823 & 0.1372 & 0.006\\
\bottomrule
\end{tabular}
\end{table}

LOF is clearly strongest on this experiment. PhiGraph is second by PR-AUC and
precision at budget, but its ROC-AUC is slightly below Isolation Forest.
This single fixed experiment establishes only that the implementation produces
a non-degenerate ranking under the reported protocol. It does not support a
claim of general or state-of-the-art superiority.

\section{Threats to Validity and Limitations}

\paragraph{Dataset validity.}
CIC-IDS2017 has documented problems in flow construction, feature generation,
and labels \cite{liu2022errors,lanvin2022errors}. Results may change on
corrected data. The machine-learning CSV also omits source and destination IP
addresses, timestamps, users, devices, and processes. The evaluated graph is
therefore a feature-similarity graph rather than the heterogeneous operational
graph for which PhiGraph is designed.

\paragraph{Experimental validity.}
The external result uses one file, one attack family, one seed, and one alert
budget. It reports no uncertainty intervals, ablation confidence intervals, or
prospective evaluation. The benchmark is balanced and does not reproduce a
production base rate. The PhiGraph score is batch-transductive because its
component percentile ranks are calculated over the complete evaluation batch;
scores can therefore change with batch composition and are not calibrated
online anomaly scores. Additional chronological splits, repeated seeds,
corrected datasets, heterogeneous event sources, and cost-sensitive metrics
are required.

\paragraph{Software validity.}
Automated tests can encode incorrect assumptions and do not substitute for
penetration testing, formal verification, independent replication, or
operational incident analysis. The ledger's hash chain detects accidental or
unauthorized changes only under an appropriate trust and key-management model.
JSON and SQLite are evaluation backends; production use requires hardened
storage, secrets management, authorization, monitoring, and backup procedures.

\paragraph{Causal and governance claims.}
Graph connectivity, spectral structure, and anomaly scores do not identify
causal effects without assumptions and suitable interventions. PhiGraph records
candidate causal pathways but does not prove causality from observational data.
Likewise, a recorded approval may still be mistaken or malicious. The platform
improves inspectability and enforcement points; it does not guarantee ethical,
legal, or correct outcomes.

\section{Reproducibility and Availability}

PhiGraph Core 4.0.0 is implemented in Python and distributed under the MIT
License. Its development repository is maintained at
\url{https://github.com/wcalmels/phigraph}. A versioned public archive should be
created together with the Zenodo record. The paper source and references are
licensed under CC BY 4.0.

The software test suite can be reproduced with:

\begin{verbatim}
python -m venv .venv
pip install -e ".[api,benchmark,dev,auth,app]"
pytest
\end{verbatim}

The CIC-IDS2017 experiment requires the Friday DDoS CSV and is run with:

\begin{verbatim}
python scripts/validate_cicids2017.py /path/to/Friday...DDos.csv
\end{verbatim}

The committed validation artifact records metrics and predictions but does not
record the source-dataset checksum, exact dependency versions, hardware,
operating system, or Git revision. It is therefore an auditable result artifact,
not yet an exact reproducibility package. A final archive should preserve the
dataset provenance and checksum, Python and dependency versions, seed,
configuration, hardware description, and Git revision.
The Zenodo DOI should be added to this section and to the repository citation
metadata after the archival deposit is published.

\section{Ethical and Operational Considerations}

PhiGraph can store sensitive evidence, identity metadata, and operational
proposals. Deployers must minimize collection, define retention periods,
encrypt transport and storage, separate tenant access, protect signing keys,
and audit privileged operations. Security telemetry and network datasets may
contain personal or confidential information and require an appropriate legal
basis. The system should be evaluated in replay or shadow mode before any
higher-authority integration. Human approval should remain mandatory where
consequences are difficult to reverse or affect safety, rights, employment,
finance, health, or critical infrastructure.

\section{Conclusion}

PhiGraph makes a specific architectural choice: intelligence does not imply
authority. Claims must remain linked to evidence and verification; actions must
remain subject to policy and explicit execution boundaries; outcomes must
remain auditable. Version 4.0.0 demonstrates this design through a typed
protocol, tamper-evident ledger, policy-gated runtime, scope-tagged records, and
an automated test suite. The external CIC-IDS2017 experiment shows a functional
relational anomaly score but also illustrates why bounded interpretation is
necessary: a conventional LOF baseline performs best, and the dataset cannot
exercise PhiGraph's full heterogeneous graph model. Future work should
prioritize independent replication, corrected and chronological security
datasets, multi-source operational graphs, policy-adversarial testing,
cryptographically anchored audit storage, and prospective shadow deployments.

\section*{Author Contributions}

Walter Calmels von Dem Knesebeck conceived the system, implemented the
software, designed the evaluation, analyzed the results, and wrote the
manuscript.

\section*{Conflict of Interest}

The author is affiliated with TUCH Systems, the organization developing
PhiGraph. This relationship may create a commercial interest in the system.
The evaluation and limitations are reported to enable independent scrutiny.

\section*{License}

This paper is licensed under the Creative Commons Attribution 4.0
International License (CC BY 4.0).

\bibliographystyle{plain}
\bibliography{references}

\end{document}
